\subsection{\problemblock{*Earth} Data Block}\label{sec:block-earth}

TAOS requires an earth model to generate a trajectory. An earth model consists of
an earth or planetary shape, a gravity model, and additional parameters such as the
rotation rate. Four standard families are provided: spherical earth, WGS-72,
WGS-84, and GEM-T1. If none is suitable, individual parameters can be overridden
to form a user-defined model.

The earth is represented as an ellipsoid whose equatorial radius is greater than its
polar radius. Shape is specified by equatorial radius and one of polar radius,
eccentricity, or flattening; once one of those three is known, the other two can be
computed. Gravity is specified by the gravity constant $GM$ and either zonal
harmonics $J_n$ or the full coefficients $C_{n,m}$ and $S_{n,m}$. Rotation rate
and the pounds-mass-to-slug conversion factor complete the model.

\taossubhead{Spherical Earth Model}

The simplest model is spherical. It requires radius, rotation rate, gravity constant,
and the lbm-to-slug conversion. Defaults are based on WGS-84 and are listed in
\cref{tab:spherical-earth-model}. The block

\begin{lstlisting}[style=taosmanual]
*earth spherical omega=0
\end{lstlisting}

selects the spherical defaults and then overrides the rotation rate.

\begin{table}[htbp]
  \centering
  \caption{Spherical earth model.}
  \label{tab:spherical-earth-model}
  \begin{tabularx}{0.76\linewidth}{@{}>{\ttfamily}lXl@{}}
    \toprule
    \normalfont Symbol & Description & Default value \\
    \midrule
    reqtr & Earth radius (ft) & 20925646.3255 \\
    omega & Rotation rate (rad/s) & $7.292115\times10^{-5}$ \\
    gm    & Gravity constant (ft\textsuperscript{3}/s\textsuperscript{2}) & $1.40764438125\times10^{16}$ \\
    g     & lbm-to-slug conversion & 32.1740485 \\
    \bottomrule
  \end{tabularx}
\end{table}

\taossubhead{WGS-72 and WGS-84 Earth Models}

The \texttt{wgs-72} and \texttt{wgs-84} models use an ellipsoidal earth. One of
those keywords must immediately follow \problemblock{*earth}, and any following
assignments override the defaults. For example,

\begin{lstlisting}[style=taosmanual]
*earth wgs-84 omega=0
\end{lstlisting}

selects WGS-84 but makes the earth nonrotating, a simplification often used in
conceptual-design studies.

Because the model is ellipsoidal, additional parameters are required as shown in
\cref{tab:wgs-earth-model-values}. Shape is defined by equatorial radius and one
of polar radius, eccentricity, or flattening. Supplying more than one of the latter
usually creates redundant and potentially inconsistent geometry.

\begin{table}[htbp]
  \centering
  \caption{WGS-72 and WGS-84 earth-model values.}
  \label{tab:wgs-earth-model-values}
  \small
  \begin{tabularx}{0.98\linewidth}{@{}>{\ttfamily}lXll@{}}
    \toprule
    \normalfont Symbol & Description & WGS-84 & WGS-72 \\
    \midrule
    reqtr & Equatorial radius (ft) & 20925646.3255 & 20925639.7638 \\
    rpolr & Polar radius (ft) & 20855486.5953 & 20855480.7087 \\
    ecc & Eccentricity & 0.0818191908426 & 0.0818188106627 \\
    flat & Flattening parameter & $1/298.257223563$ & $1/298.26$ \\
    omega & Rotation rate (rad/s) & $7.292115\times10^{-5}$ & $7.292115147\times10^{-5}$ \\
    g & lbm-to-slug conversion & 32.1740485 & 32.1740485 \\
    gm & Gravity constant (ft\textsuperscript{3}/s\textsuperscript{2}) & $1.40764438125\times10^{16}$ & $1.40764544069\times10^{16}$ \\
    j2 & Degree-2 zonal harmonic & $1.08262998905\times10^{-3}$ & $1.08261579002\times10^{-3}$ \\
    \bottomrule
  \end{tabularx}
\end{table}

The standard WGS gravity model uses the second-degree zonal harmonic $J_2$ to
account for oblateness; higher coefficients are neglected.

\taossubhead{TSAP-Compatible Earth Models}

TSAP and PMAST used WGS-72 and WGS-84 models containing zonal harmonics
through degree 4. $J_2$ represents earth oblateness, while $J_3$ and $J_4$ describe
additional latitude-dependent variations. Their use is questionable for typical
trajectories because they are comparable to omitted nonzonal harmonics and to the
neglected gravitational effects of the sun and moon. Nevertheless, the
\texttt{tsap-72} and \texttt{tsap-84} models retain them so TAOS results can be
consistent with its predecessors.

\begin{table}[htbp]
  \centering
  \caption{TSAP-compatible earth-model values.}
  \label{tab:tsap-compatible-earth-values}
  \begin{tabularx}{0.84\linewidth}{@{}>{\ttfamily}lXll@{}}
    \toprule
    \normalfont Symbol & Description & WGS-84 & WGS-72 \\
    \midrule
    j3 & Degree-3 zonal harmonic & $-2.532153068\times10^{-6}$ & $-2.538810043\times10^{-6}$ \\
    j4 & Degree-4 zonal harmonic & $-1.610987610\times10^{-6}$ & $-1.655970000\times10^{-6}$ \\
    \bottomrule
  \end{tabularx}
\end{table}

\taossubhead{Full Degree-4 Earth Models}

Two additional models, \texttt{wgs-84-full} and \texttt{gem-t1-full}, include all
geopotential terms through degree 4. Their input coefficients are unnormalized, so
replacement values must also be unnormalized. The source does not recommend
these models for high-accuracy work because neglected solar and lunar gravity is of
the same order as many retained coefficients.

\addtocounter{table}{1}% The source manual intentionally skips Table 4-18.
\begin{longtable}{@{}>{\ttfamily}lp{0.33\linewidth}ll@{}}
  \caption{Full WGS-84 and GEM-T1 earth-model values.}
  \label{tab:full-earth-model-values}\\
  \toprule
  \normalfont Symbol & Description & WGS-84 & GEM-T1 \\
  \midrule
  \endfirsthead
  \toprule
  \normalfont Symbol & Description & WGS-84 & GEM-T1 \\
  \midrule
  \endhead
  reqtr & Equatorial radius (ft) & 20925646.3255 & 20925646.3255 \\
  rpolr & Polar radius (ft) & 20855486.5953 & 20855486.5953 \\
  ecc & Eccentricity & 0.0818191908426 & 0.0818191908426 \\
  flat & Flattening parameter & $1/298.257223563$ & $1/298.257223563$ \\
  omega & Rotation rate (rad/s) & $7.292115\times10^{-5}$ & $7.292115\times10^{-5}$ \\
  g & lbm-to-slug conversion & 32.1740485 & 32.1740485 \\
  gm & Gravity constant & $1.40764438125\times10^{16}$ & $1.4076441552\times10^{16}$ \\
  c20 & Degree-2 zonal harmonic & $-1.082629\times10^{-3}$ & $-1.082625\times10^{-3}$ \\
  c22 & Degree-2 nonzonal harmonic & $1.572805\times10^{-6}$ & $1.574322\times10^{-6}$ \\
  c30 & Degree-3 zonal harmonic & $2.532153\times10^{-6}$ & $2.532618\times10^{-6}$ \\
  c31 & Degree-3 nonzonal harmonic & $2.194673\times10^{-6}$ & $2.192402\times10^{-6}$ \\
  c32 & Degree-3 nonzonal harmonic & $3.096837\times10^{-7}$ & $3.086210\times10^{-7}$ \\
  c33 & Degree-3 nonzonal harmonic & $1.000789\times10^{-7}$ & $1.005372\times10^{-7}$ \\
  c40 & Degree-4 zonal harmonic & $1.610987\times10^{-6}$ & $1.616190\times10^{-6}$ \\
  c41 & Degree-4 nonzonal harmonic & $-5.080013\times10^{-7}$ & $-5.060561\times10^{-7}$ \\
  c42 & Degree-4 nonzonal harmonic & $7.780961\times10^{-8}$ & $7.759155\times10^{-8}$ \\
  c43 & Degree-4 nonzonal harmonic & $5.926679\times10^{-8}$ & $5.922238\times10^{-8}$ \\
  c44 & Degree-4 nonzonal harmonic & $-3.948164\times10^{-9}$ & $-4.015116\times10^{-9}$ \\
  s22 & Degree-2 nonzonal harmonic & $-9.023759\times10^{-7}$ & $-9.035928\times10^{-7}$ \\
  s31 & Degree-3 nonzonal harmonic & $2.709571\times10^{-7}$ & $2.695880\times10^{-7}$ \\
  s32 & Degree-3 nonzonal harmonic & $-2.121201\times10^{-7}$ & $-2.119137\times10^{-7}$ \\
  s33 & Degree-3 nonzonal harmonic & $1.973456\times10^{-7}$ & $1.970571\times10^{-7}$ \\
  s41 & Degree-4 nonzonal harmonic & $-4.498693\times10^{-7}$ & $-4.507384\times10^{-7}$ \\
  s42 & Degree-4 nonzonal harmonic & $1.466394\times10^{-7}$ & $1.484816\times10^{-7}$ \\
  s43 & Degree-4 nonzonal harmonic & $-1.189998\times10^{-8}$ & $-1.198933\times10^{-8}$ \\
  s44 & Degree-4 nonzonal harmonic & $6.540039\times10^{-9}$ & $6.517407\times10^{-9}$ \\
  \bottomrule
\end{longtable}
